\section{Limitations and Version Scope}\label{sec:limits}

The v1 evidence comes from controlled two-dimensional point-kinematic simulators, scalar monitored margins, fixed scenario grids, Gaussian measurement noise, and one braking or serialization recovery family. It does not cover vehicle dynamics, perception failures, localization, networking, asynchronous execution, or a full navigation stack. The analytical stopping deadline is deliberately conservative, while the replay-derived $8.8$\,s deadline is empirical for one deterministic episode and a $0.1$\,s trigger grid.

The dose-response and boundary figures are exploratory: fixed-pad selection and reporting use the same runs, and the predictive horizon remains at $2.0$\,s. The second single-margin experiment separates tuning and test seeds and gives both methods the same candidate grid, but only noise realizations are held out. The simulator, factor levels, and selection rule are fixed rather than sampled from a target population. Its reported Wilson intervals are pointwise, and their largest upper endpoint is not a simultaneous confidence bound. Because the held-out primary arms have different collision outcomes, their conditional completion times are descriptive and no inferential efficiency comparison is made.

In the coordination simulator, untreated conflict trajectories physically overlap, and waiting robots have explicit blockers and release conditions. The first trigger nevertheless commits an immutable ordering; the simulator does not re-plan. The observed timing advantage therefore concerns this one-shot serialization rather than multi-robot coordination in general. The bootstrap interval varies measurement noise over a fixed 27-scenario grid and does not cover new geometries, dynamics, robot counts, or degradation families. The no-intervention arm consumes no measurement noise, so repeated runs within a physical scenario are deterministic duplicates rather than independent safety evidence. The model contains only three robots and one shared zone.

ROS~2/Nav2 experiments with localization degradation, communication faults, planner interactions, and full-stack telemetry are reserved for v2; no v1 claim extends to them. The related-work search is targeted rather than systematic. A later archival submission should broaden and document the search and recheck each citation against its surrounding claim.
